\documentclass[tikz,border=10pt]{standalone}
\usetikzlibrary{shapes.geometric, arrows, positioning}
\usepackage{tikz}
\usepackage{ctex}
\usetikzlibrary{decorations.pathreplacing}
\usetikzlibrary{patterns}
\tikzstyle{block} = [rectangle, rounded corners, minimum width=2.5cm, minimum height=1cm,text centered, draw=black, fill=blue!20]
\tikzstyle{attention_block} = [rectangle, rounded corners, minimum width=2.5cm, minimum height=1cm,text centered, draw=black, fill=green!20]
\tikzstyle{convblock} = [rectangle, rounded corners, minimum width=2.5cm, minimum height=1cm,text centered, draw=black, fill=orange!20]
\tikzstyle{arrow} = [thick,->,>=stealth]

\begin{document}
\begin{tikzpicture}
    % Định nghĩa chiều rộng và chiều cao của ô mảng
    \def\w{0.8}
    \def\h{0.8}

    % Vẽ tất cả các ô và đường viền, gồm cả phần dấu chấm lửng
    \foreach \x in {0,...,10} {
        \ifnum \x<4
            \fill[fill=green!30] (\x*\w,0) rectangle ++ (\w,\h);
        \else \ifnum \x=4
            \fill[fill=yellow!30] (\x*\w,0) rectangle ++ (\w,\h);
        \else \ifnum \x=5
            % Ô dấu chấm lửng
            \fill[fill=white] (\x*\w,0) rectangle ++ (\w,\h);
        \else \ifnum \x<10
            \fill[fill=green!30] (\x*\w,0) rectangle ++ (\w,\h);
        \else
            \fill[fill=yellow!30] (\x*\w,0) rectangle ++ (\w,\h);
        \fi\fi\fi\fi
        \draw (\x*\w,0) rectangle ++ (\w,\h);
    }

    % Thêm chữ cái
    \foreach \pos/\letter/\col in {0/a/black,1/c/black,2/a/black,3/c/black,4/e/red} {
        \node[anchor=center] at (\pos*\w+0.4,0.4) {\textcolor{\col}{\letter}};
    }

    \node at (5*\w+0.4,0.4) {$\cdots$};

    \foreach \pos/\letter/\col in {6/a/black,7/c/black,8/a/black,9/c/black,10/c/red} {
        \node[anchor=center] at (\pos*\w+0.4,0.4) {\textcolor{\col}{\letter}};
    }

    % Thêm ngoặc móc và chú thích
    \draw[decoration={brace},decorate] (0,\h+0.1) -- (3.2,\h+0.1)
        node[midway,above=6pt] {$s[0 \ldots \pi[i]-1]$};

    \draw[decoration={brace},decorate] (4.8,\h+0.1) -- (8.0,\h+0.1)
        node[midway,above=6pt] {$s[i-\pi[i]-1 \ldots i]$};

    % Thêm chú thích mũi tên
    \draw[->] (3.6,\h+0.1) -- (3.6,\h+0.8);
    \node[above] at (3.6,\h+0.8) {$s[\pi[i]]$};

    \draw[decorate,decoration={brace,mirror,amplitude=6pt}] (0,-0.2) -- (8.0,-0.2)
        node[midway,below=8pt] {$s[0 \ldots i]$};

    \draw[->] (8.4,\h+0.1) -- (8.4,\h+0.8);
    \node[above] at (8.4,\h+0.8) {$s[i+1]$};

    % Thêm chú thích không khớp ở ngoài cùng bên phải
    \node[right] at (11*\w,0.4) {$s[\pi[i]] \neq s[i+1]$ không khớp};
\end{tikzpicture}

\begin{tikzpicture}
    % Định nghĩa chiều rộng và chiều cao của ô mảng
    \def\w{0.8}
    \def\h{0.8}

    % Vẽ tất cả các ô và đường viền, gồm cả phần dấu chấm lửng
    \foreach \x in {0,...,10} {
        \ifnum \x<2
            \fill[fill=green!30] (\x*\w,0) rectangle ++ (\w,\h);
        \else \ifnum \x=2
            \fill[fill=yellow!30] (\x*\w,0) rectangle ++ (\w,\h);
        \else \ifnum \x<5
            \fill[fill=gray!20] (\x*\w,0) rectangle ++ (\w,\h);
        \else \ifnum \x=5
            % Ô dấu chấm lửng
            \fill[fill=gray!20] (\x*\w,0) rectangle ++ (\w,\h);
        \else \ifnum \x<8
            \fill[fill=gray!20] (\x*\w,0) rectangle ++ (\w,\h);
        \else \ifnum \x<10
            \fill[fill=green!30] (\x*\w,0) rectangle ++ (\w,\h);
        \else
            \fill[fill=yellow!30] (\x*\w,0) rectangle ++ (\w,\h);
        \fi\fi\fi\fi\fi\fi
        \draw (\x*\w,0) rectangle ++ (\w,\h);
    }

    % Thêm chữ cái
    \foreach \pos/\letter/\col in {0/a/black,1/c/black,2/a/red,3/c/black,4/e/black} {
        \node[anchor=center] at (\pos*\w+0.4,0.4) {\textcolor{\col}{\letter}};
    }

    \node at (5*\w+0.4,0.4) {$\cdots$};

    \foreach \pos/\letter/\col in {6/a/black,7/c/black,8/a/black,9/c/black,10/c/red} {
        \node[anchor=center] at (\pos*\w+0.4,0.4) {\textcolor{\col}{\letter}};
    }

    % Thêm chú thích mũi tên
    \draw[->] (2*\w+0.4,\h+0.1) -- (2*\w+0.4,\h+0.8);
    \node[above] at (2*\w+0.4,\h+0.8) {$s[j]$};

    \draw[->] (6.8,\h+0.1) -- (6.8,\h+0.3);
    \node[above] at (6.8,\h+0.3) {$s[i-j+1]$};

    \draw[->] (10*\w+0.4,\h+0.1) -- (10*\w+0.4,\h+0.8);
    \node[above] at (10*\w+0.4,\h+0.8) {$s[i+1]$};

    % Thêm chú thích ngoặc móc phía dưới
    \draw[decorate,decoration={brace,mirror,amplitude=6pt}] (0,-0.2) -- (1.6,-0.2)
        node[midway,below=8pt] {Độ dài thứ hai $j$};

    \draw[decorate,decoration={brace,mirror,amplitude=6pt}] (8*\w,-0.2) -- (8.0,-0.2)
        node[midway,below=8pt] {Độ dài thứ hai $j$};

    % Thêm dòng chữ ngoài cùng bên phải
    \node[right] at (11*\w,0.4) {Tìm kết quả khớp có thể};
\end{tikzpicture}

\begin{tikzpicture}
    % Định nghĩa chiều rộng và chiều cao của ô mảng
    \def\w{0.8}
    \def\h{0.8}

    % Vẽ mảng phía trên
    \foreach \x in {0,...,10} {
        \ifnum \x<4
            \fill[fill=green!30] (\x*\w,0) rectangle ++ (\w,\h);
        \else \ifnum \x=4
            \fill[fill=white] (\x*\w,0) rectangle ++ (\w,\h);
        \else \ifnum \x=5
            \fill[fill=white] (\x*\w,0) rectangle ++ (\w,\h);
        \else \ifnum \x<8
            \fill[fill=gray!20] (\x*\w,0) rectangle ++ (\w,\h);
        \else \ifnum \x<10
            \fill[fill=green!30] (\x*\w,0) rectangle ++ (\w,\h);
        \else
            \fill[fill=white] (\x*\w,0) rectangle ++ (\w,\h);
        \fi\fi\fi\fi\fi
        \draw (\x*\w,0) rectangle ++ (\w,\h);
    }

    % Thêm chữ cái
    \foreach \pos/\letter in {0/a,1/c,2/a,3/c,4/e} {
        \node[anchor=center] at (\pos*\w+0.4,0.4) {\letter};
    }

    \node at (5*\w+0.4,0.4) {$\cdots$};

    \foreach \pos/\letter in {6/a,7/c,8/a,9/c,10/c} {
        \node[anchor=center] at (\pos*\w+0.4,0.4) {\letter};
    }

    % Vẽ mảng phía dưới
    \foreach \x in {0,...,3} {
        \fill[fill=green!30] (\x*\w,-2.6) rectangle ++ (\w,\h);
        \draw (\x*\w,-2.6) rectangle ++ (\w,\h);
        \ifcase\x
            \node[anchor=center] at (\x*\w+0.4,-2.2) {a};
        \or
            \node[anchor=center] at (\x*\w+0.4,-2.2) {c};
        \or
            \node[anchor=center] at (\x*\w+0.4,-2.2) {a};
        \or
            \node[anchor=center] at (\x*\w+0.4,-2.2) {c};
        \fi
    }

    % Thêm chú thích ngoặc móc phía trên
    \draw[decoration={brace},decorate] (0,\h+0.1) -- (3.2,\h+0.1)
        node[midway,above=6pt] {$s[0 \ldots \pi[i]-1]$};

    \draw[decoration={brace},decorate] (4.8,\h+0.1) -- (8.0,\h+0.1)
        node[midway,above=6pt] {$s[i-\pi[i]+1 \ldots i]$};

    % Thêm chú thích hình elip
    \draw[blue, thick] (0.8,0.4) ellipse (0.8cm and 0.4cm);
    \draw[blue, thick] (7.2,0.4) ellipse (0.8cm and 0.4cm);

    % Thêm chú thích phía dưới
    \draw[decoration={brace,mirror},decorate] (6.4,-0.1) -- (8.0,-0.1)
        node[midway,below=6pt] {$s[i-j+1 \ldots i]$};

    \draw[decoration={brace,mirror},decorate] (0,-2.8) -- (1.6,-2.8)
        node[midway,below=6pt] {$s[0 \ldots j-1]$};

    \draw[decoration={brace},decorate] (1.6,-1.6) -- (3.2,-1.6)
        node[midway,above=6pt] {$s[\pi[i]-j \ldots \pi[i]-1]$};

    % Thêm elip phía dưới
    \draw[blue, thick] (0.8,-2.2) ellipse (0.8cm and 0.4cm);
    \draw[blue, thick] (2.4,-2.2) ellipse (0.8cm and 0.4cm);

    % Thêm liên kết nét đứt
    \draw[gray, dashed, ->] (7.2,0) -- (2.4,-1.6);
    \draw[gray, dashed, ->] (0.8,0) -- (0.8,-1.6);

    \draw[->] (3.5*\w,-2.8) -- (3.5*\w,-3.6);
    \node[below] at (3.5*\w,-3.6) {$s[\pi[i]-1]$};
    \node[below] at (3.5*\w,-4.2) {Giá trị hàm tiền tố $\pi[\pi[i]-1]=2=j$};

    % Thêm phần giải thích bên phải
    \node[right] at (6*\w,-2.2) {Độ dài thứ hai $j$ tương đương với xâu con};
    \node[right] at (6*\w,-2.8) {Giá trị hàm tiền tố của $s[0 \ldots \pi[i]-1]$};
\end{tikzpicture}


\begin{tikzpicture}
    % Định nghĩa chiều rộng và chiều cao của ô
    \def\w{0.8}
    \def\h{0.8}

    % Vẽ phần chữ cái phía trên
    \foreach \x/\letter in {
        0/a,1/b,2/a,3/b,4/c,5/\#,6/a,7/b,8/c,9/d,10/e,11/f,12/a,13/b,14/a,15/b,
        16/g,17/a,18/b,19/a,20/b,21/c,22/f,23/g
    } {
        \ifnum \x<5
            \fill[fill=green!20] (\x*\w,0) rectangle ++ (\w,\h);
        \else \ifnum \x=5
            \fill[fill=gray!30] (\x*\w,0) rectangle ++ (\w,\h);
        \else \ifnum \x>16 \ifnum \x<21
            \fill[fill=green!20] (\x*\w,0) rectangle ++ (\w,\h);
        \else \ifnum \x=21
            \fill[fill=green!30] (\x*\w,0) rectangle ++ (\w,\h);
        \fi\fi\fi\fi\fi
        \draw (\x*\w,0) rectangle ++ (\w,\h);
        \node[anchor=center] at (\x*\w+0.4,0.4) {\letter};
    }

    % Vẽ phần số phía dưới
    \foreach \x/\num in {
        0/0,1/0,2/1,3/2,4/0,5/0,6/1,7/2,8/0,9/0,10/0,11/0,12/1,13/2,14/3,15/4,
        16/0,17/1,18/2,19/3,20/4,21/5,22/0,23/0
    } {
        \draw (\x*\w,-\h) rectangle ++ (\w,\h);
        \node[anchor=center] at (\x*\w+0.4,-0.4) {\num};
    }

    % Thêm chú thích và mũi tên phía dưới
    % Chú thích n
    \draw[decoration={brace,mirror},decorate] (0,-1.2) -- (4.0,-1.2);
    \node[below] at (2,-1.3) {Xâu s có độ dài $n$};

    % Chú thích n (bên phải)
    \draw[decoration={brace,mirror},decorate] (17*\w,-1.2) -- (22*\w,-1.2);
    \node[below] at (19.5*\w,-1.3) {Xâu con khớp độ dài $n$};

    % Chú thích mũi tên dọc
    \draw[->] (6.5*\w,-1.0) -- (6.5*\w,-3.2);
    \node[below] at (6.5*\w,-3.2) {Chỉ số $[n+1]$};

    \draw[->] (17.5*\w,-1.0) -- (17.5*\w,-3.2);
    \node[below] at (17.5*\w,-3.2) {Chỉ số $[i-(n-1)]$};

    \draw[->] (21.5*\w,-1.0) -- (21.5*\w,-2.4);
    \node[below] at (21.5*\w,-2.4) {Chỉ số $[i]$ thỏa mãn $\pi[{i}]=n$};

    % Mũi tên ngang
    \draw[->, thick] (6.5*\w,-2.5) -- (17.5*\w,-2.5);
    \node[below] at (11.5*\w,-2.5) {Chỉ số của xâu con trong văn bản t là $i-2n$};

\end{tikzpicture}

\end{document}
